\subsection{LOCALLY CONNECTED SPACES}

DEFINITION 4. A topological space X is said to be locally connected if each point of X has a fundamental system of connected neighbourhoods.

* We shall see in Chapter IV, § 2, no. 5, that the real line is a locally connected space. *

The existence, at each point x of a space X, of one connected neighbourhood of x by no means implies that X is locally connected. In particular, X can be connected but not locally connected (Exercises 2 and 13). Conversely, a space can be locally connected but not connected (e.g.~a discrete space which contains more than one point).

PROPOSITION 11. A necessary and sufficient condition for a space X to be locally connected is that every component of an open set in X is open in X.

The condition is sufficient, since the component of x relative to an open neighbourhood of x is then a neighbourhood of x in X.

Conversely, let A be an open subset of a locally connected space X, let B be a component of A, and let \(x \in B\) . Let V be a connected neighbourhood of x contained in A; by the definition of components, V is contained in B; hence B is open in X (§ 1, no. 2, Proposition 1).

The components of a locally connected space X therefore form a partition of X into open sets, and hence X is the sum (§ 2, no. 4, Example 3) of its components.

COROLLARY. Let U be an open subset of a locally connected space X, and let V be a component of U. Then the frontier of V (relative to X) is contained in the frontier of U.

For V is open and closed in U, hence a frontier point of V (relative to X) cannot belong to U, for it would also be a frontier point of V relative to U, and there is none.

PROPOSITION 12. Every quotient space of a locally connected space is locally connected.

Let X be a locally connected space, R an equivalence relation on X, \(\varphi: X \to X/R\) the canonical mapping. Let A be an open subset of

X/R and C a component of A. Then \(\overline{\varphi}^{1}(C)\) is a union of components of \(\overline{\varphi}^{1}(A)\) ; for if \(x \in \overline{\varphi}^{1}(C)\) and if K is the component of \(x\) in \(\overline{\varphi}^{1}(A)\) then \(\varphi(K)\) is connected (no. 2, Proposition 4), is contained in A, and contains \(\varphi(x)\) ; hence \(\varphi(K) \subset C\) by the definition of C, and therefore \(K \subset \overline{\varphi}^{1}(C)\) . Since X is locally connected and \(\overline{\varphi}^{1}(A)\) is open in X, it follows from Proposition 11 that \(\overline{\varphi}^{1}(C)\) is open in X; consequently C is open in X/R and hence, by Proposition 11 again, X/R is locally connected.

PROPOSITION 13. a) Let \((\mathbf{X}_t)_{t \in I}\) be a family of locally connected spaces such that \(\mathbf{X}_t\) is connected for all but a finite number of indices \(t \in I\) . Then the product space \(\mathbf{X} = \prod_{i \in I} \mathbf{X}_t\) is locally connected.

\begin{enumerate}
\def\labelenumi{\alph{enumi})}
\setcounter{enumi}{1}
\item
  Conversely, if the product of a family \((\mathbf{X}_{t})\) of non-empty topological spaces is locally connected, then each \(X_{t}\) is locally connected, and \(X_{t}\) is connected for all but a finite number of indices.
\item
  Let \(\mathbf{J}\) be the finite subset of \(\mathbf{I}\) such that \(\mathbf{X}_i\) is not connected if and only if \(i \in \mathbf{J}\) . Let
\end{enumerate}

\[
\mathrm{U} = \prod_ {i \in I} \mathrm{U} _ {i}
\]

be an elementary set containing a point \(x = (x_{t})\) of \(\mathbf{X}\) and let \(\mathbf{K}\) be the finite subset of \(\mathbf{I}\) such that \(\mathbf{U}_{t} \neq \mathbf{X}_{t}\) if and only if \(t \in \mathbf{K}\) . Let \(\mathbf{V}_{t}\) be \(\mathbf{X}_{t}\) for \(t \notin J \cup K\) , and let \(V_{t}\) be a connected neighbourhood of \(x_{t}\) contained in \(U_{t}\) for \(t \in J \cup K\) ; then

\[
\mathbf {V} = \prod_ {i \in I} \mathbf {V} _ {i}
\]

is connected (by Proposition 8 of no. 4) and is a neighbourhood of x contained in U. Hence X is locally connected.

\begin{enumerate}
\def\labelenumi{\alph{enumi})}
\setcounter{enumi}{1}
\tightlist
\item
  Let \(a = (a_{i})\) be a point of \(\mathbf{X}\) and let \(\mathbf{V}\) be a connected neighbourhood of \(a\) in \(\mathbf{X}\) . Since we have \(\mathrm{pr}_{\mathfrak{i}}\mathbf{V} = \mathbf{X}_{\mathfrak{i}}\) except for a finite number of indices (§ 4, no. 1) it follows from no. 2, Proposition 4 that the \(\mathbf{X}_{\mathfrak{i}}\) are connected, for all but a finite number of indices. On the other hand, for each \(x \in I\) , each \(a_{x} \in X_{x}\) and each neighbourhood \(V_{x}\) of \(a_{x}\) in \(X_{x}\) , there is a point \(x\) of \(\mathbf{X}\) such that \(\mathrm{pr}_{\mathfrak{x}}x = a_{x}\) , and
\end{enumerate}

\[
\mathbf {V} = \mathbf {V} _ {x} \times \prod_ {i \neq x} \mathbf {X} _ {i}
\]

is a neighbourhood of \(x\) in \(\mathbf{X}\) ; \(\mathbf{V}\) therefore contains a connected neighbourhood \(W\) of \(x\) , whose projection \(\mathrm{pr}_xW\) is a connected neighbourhood of \(a_x\) contained in \(V_x\) (no. 2, Proposition 4 and § 4, no. 2, Proposition 5). Hence each \(X_x\) is locally connected.

